\section{Data Appendix}\label{app:data}

\subsection{Fact 1: Volatility of Expectations (Dividend Growth)}\label{app:data_fact1_expectations_volatility}

\paragraph{Overview.}
This subsection describes the replication of Fact~1 in Section~\ref{sec:facts.sec} of the paper: the variance share of (i) an ``objective'' dividend-growth forecast constructed from dividend dynamics and (ii) a survey-based ``subjective'' dividend-growth forecast.

\paragraph{Inputs.}
The replication uses two quarterly time series from the dividend-growth expectations dataset:
\begin{itemize}
\item \textit{Realized next-year dividend growth} \(\Delta d_{t,t+4}\) (realized next-year log dividend growth at quarterly dates);
\item \textit{Subjective expectations} \(\mathbb{E}^{sub}[\Delta d_{t,t+4}]\) (I/B/E/S survey expectations of next-year log dividend growth).
\end{itemize}
To construct the objective forecast, we use the Shiller historical dataset and its real dividend series to build quarterly dividend growth.

\paragraph{Sample.}
The baseline replication uses a quarterly sample spanning 2003Q1--2021Q3 (75 observations), dictated by the overlap between the realized next-year dividend-growth series and the survey expectations series. The variance ratios are computed on the common non-missing sample for each expectations measure.

\paragraph{Objective expectations construction.}
Let \(g^d_t\) denote quarterly log dividend growth constructed from Shiller real dividends by taking the end-of-quarter (last monthly) observation in each quarter and forming log differences. We estimate an AR(1) on \(g^d_t\):
\[
g^d_t = a + \rho g^d_{t-1} + \varepsilon_t.
\]
Given the estimated \((a,\rho)\), it maps quarterly growth into an implied one-year-ahead expectation by iterating the AR(1) forward and summing the implied quarterly growth rates:
\[
\mathbb{E}^{obj}[\Delta d_{t,t+4}] \equiv \sum_{h=1}^{4}\mathbb{E}_t[g^d_{t+h}],
\qquad
\mathbb{E}_t[g^d_{t+h}] = \mu + \rho^h\,(g^d_t-\mu),\ \mu=\frac{a}{1-\rho}.
\]
The resulting objective series is then aligned to the quarterly dates in the dividend-growth expectations dataset by matching year/quarter identifiers.

\paragraph{Variance ratios.}
For each expectations measure \(m\in\{obj,sub\}\), the replication computes:
\[
\frac{\mathrm{Var}\!\bigl[\mathbb{E}^{m}[\Delta d_{t,t+4}]\bigr]}{\mathrm{Var}[\Delta d_{t,t+4}]},
\]
using the common sample for which both the realized series \(\Delta d_{t,t+4}\) and the corresponding expectation series are non-missing (objective and subjective samples coincide in the default run).

\paragraph{Outputs and replicated numbers.}
We report the variance ratios implied by the baseline inputs used for the paper. The dividend-growth variance ratios are:
\[
\frac{\mathrm{Var}\!\bigl[\mathbb{E}^{obj}[\Delta d_{t,t+4}]\bigr]}{\mathrm{Var}[\Delta d_{t,t+4}]} \;=\; 0.0486 \;(\approx 0.05),
\qquad
\frac{\mathrm{Var}\!\bigl[\mathbb{E}^{sub}[\Delta d_{t,t+4}]\bigr]}{\mathrm{Var}[\Delta d_{t,t+4}]} \;=\; 0.8268 \;(\approx 0.83).
\]

\begin{table}[t!]
    \centering
    \caption{Objective Dividend-Growth Expectations from AR(1): Sample Window Robustness}
    \label{tab:fact1_ar1_windows}
    \begin{threeparttable}
    \begin{tabular}{lcc}
    \toprule
    Sample window & Observations ($N$) & Variance ratio \\
    \midrule
    Full (1871Q2--2022Q2) & 605 & 0.0665 \\
    Post-1947 (1947Q1--2022Q2) & 302 & 0.0599 \\
    Post-1976 (1976Q1--2022Q2) & 186 & 0.0501 \\
    Post-2003 (2003Q1--2022Q2) & 78 & 0.0495 \\
    \bottomrule
    \end{tabular}
    \begin{tablenotes}[flushleft]
    \footnotesize
    \item Notes: The objective expectation is constructed from an AR(1) estimated on quarterly log dividend growth, $\Delta d_{t+1}=a+\rho\Delta d_t+\varepsilon_{t+1}$. The one-year-ahead expectation is the implied four-quarter forecast, $E_t[\Delta d_{t,t+4}]$. The reported variance ratio is $\mathrm{Var}\!\left(E_t[\Delta d_{t,t+4}]\right)\big/\mathrm{Var}\!\left(\Delta d_{t,t+4}\right)$.
    \end{tablenotes}
    \end{threeparttable}
\end{table}


\subsection{Fact 1 Robustness: Long-Sample Objective Expectations}\label{app:data_fact1_long_sample_robustness}

\paragraph{Overview.}
This subsection reports long-sample robustness checks for the objective dividend-growth expectation in Fact~1. We keep the target object fixed,
\[
\frac{\mathrm{Var}\!\bigl[\mathbb{E}^{obj}[\Delta d_{t,t+4}]\bigr]}{\mathrm{Var}[\Delta d_{t,t+4}]},
\]
and vary only the forecasting specification and sample window.

\paragraph{Data and quarterly construction.}
All robustness exercises use the Shiller historical market series. We convert monthly data to quarterly frequency by taking end-of-quarter (last monthly) observations of real prices and real dividends, then define:
\[
\Delta d_t = \log D_t - \log D_{t-1},\qquad
\Delta d_{t,t+4} = \log D_{t+4} - \log D_t,\qquad
pd_t = \log P_t - \log D_t.
\]

\paragraph{AR(1)-implied objective expectation.}
The first robustness follows the same approach as the baseline Fact~1 construction. We estimate
\[
\Delta d_{t+1}=a+\rho \Delta d_t+\varepsilon_{t+1},
\]
form the implied four-quarter expectation \(E_t[\Delta d_{t,t+4}]\), and compute the variance ratio across alternative sample windows. Results are reported in Table~\ref{tab:fact1_ar1_windows}.

\paragraph{\textit{p/d}-only versus VAR benchmark.}
The second robustness compares two objective forecasting approaches in long samples.  
(i) A direct horizon-$4$ regression:
\[
\Delta d_{t,t+4}=a+b\,pd_t+\varepsilon_{t+4}.
\]
(ii) A quarterly VAR(1) in \((\Delta d_t,pd_t)\), from which we construct \(E_t[\Delta d_{t,t+4}]\) using model-implied four-step forecasts.  
Table~\ref{tab:fact1_pd_vs_var} reports both variance ratios for the full and post-1947 samples.

\begin{table}[t!]
    \centering
    \caption{Objective Dividend-Growth Expectations: \textit{p/d}-Only vs VAR}
    \label{tab:fact1_pd_vs_var}
    \begin{threeparttable}
    \footnotesize
    \begin{adjustbox}{max width=\textwidth}
    \begin{tabular}{lcccc}
    \toprule
    Sample window & $N$ (\textit{p/d} only) & Var. ratio (\textit{p/d} only) & $N$ (VAR) & Var. ratio (VAR) \\
    \midrule
    Full (1871Q1--2022Q2) & 606 & 0.0695 & 605 & 0.1356 \\
    Post-1947 (1947Q1--2022Q2) & 302 & 0.0020 & 302 & 0.0936 \\
    \bottomrule
    \end{tabular}
    \end{adjustbox}
    \begin{tablenotes}[flushleft]
    \item Notes: The \textit{p/d}-only specification is the direct horizon-$4$ regression, $\Delta d_{t,t+4}=a+b\,(p/d)_t+\varepsilon_{t+4}$. The VAR specification is a quarterly VAR(1) in $(\Delta d_t,p/d_t)$, and $E_t[\Delta d_{t,t+4}]$ is obtained from the model-implied four-step forecasts. In both columns, the reported statistic is $\mathrm{Var}\!\left(E_t[\Delta d_{t,t+4}]\right)\big/\mathrm{Var}\!\left(\Delta d_{t,t+4}\right)$.
    \end{tablenotes}
    \end{threeparttable}
\end{table}


\paragraph{Interpretation.}
The AR(1)-implied objective expectation remains low-variance in long samples (around 0.05--0.07), close to the baseline Fact~1 estimate. The direct \textit{p/d}-only approach contributes little in the post-1947 sample, while the VAR allows richer short-run dynamics and yields slightly larger objective-expectation variance ratios.
In all specifications, movements in expectations, as inferred by an econometrician, explain only a small fraction of the total variation in realized dividend growth.

\subsection{Fact 3: I/B/E/S Expectations and Labor (Firm-Level)}\label{app:data_fact3_expectations_labor}

\paragraph{Overview.}
This subsection describes the firm-level empirical exercise behind Fact~3 in Section~\ref{sec:facts.sec}: regressions of realized labor outcomes on lagged earnings-growth expectations from I/B/E/S with firm-clustered standard errors.

\paragraph{Data sources.}
The replication combines WRDS extracts from:
\begin{itemize}
\item Compustat Annual (\texttt{comp.funda}): \texttt{gvkey}, \texttt{datadate}, \texttt{fyear}, \texttt{emp}, \texttt{xlr};
\item CRSP monthly stock file (\texttt{crsp.msf}): \texttt{permno}, \texttt{date}, \texttt{ret};
\item CRSP--Compustat link table (\texttt{crsp.ccmxpf\_linktable});
\item I/B/E/S Summary Statistics (\texttt{ibes.statsum\_epsus});
\item WRDS I/B/E/S--CRSP historical link (\texttt{wrdsapps.ibcrsphist});
\item S\&P~500 historical membership list (local \texttt{DSP500LIST\_v4.csv}).
\end{itemize}

\paragraph{Sample selection.}
The baseline replication applies the following filters:
\begin{enumerate}
\item Compustat observations with non-missing \texttt{gvkey} and \texttt{datadate};
\item CRSP observations with valid monthly returns (\texttt{ret} $>-1$);
\item I/B/E/S observations restricted to \texttt{measure=EPS}, horizons \(\texttt{FPI}\in\{6,7,8,9\}\), and non-missing \texttt{statpers} and forecast value;
\item I/B/E/S--CRSP links with \texttt{score}\(\le 2\);
\item CRSP--Compustat linking by valid link-date intervals;
\item S\&P~500 membership filter by \texttt{permno}-date using historical entry/exit dates from \texttt{DSP500LIST\_v4.csv}.
\end{enumerate}

\paragraph{Variable construction.}
For firm \(i\) and year \(t\):
\begin{itemize}
\item Payroll outcome:
\[
Y^{pay}_{it}=\frac{\Delta xlr_{it}}{\sigma_i(\Delta xlr)},
\qquad
\Delta xlr_{it}=xlr_{it}-xlr_{i,t-1}.
\]
\item Workers outcome:
\[
Y^{emp}_{it}=\frac{\Delta emp_{it}}{\sigma_i(\Delta emp)},
\qquad
\Delta emp_{it}=emp_{it}-emp_{i,t-1}.
\]
\item Expectations regressor: from I/B/E/S, quarter-level expected one-year earnings is constructed as the sum of available EPS forecasts for \(\{6,7,8,9\}\), then transformed to firm-level annual changes and standardized:
\[
X_{it}=\frac{\Delta \widehat{earn}_{it}}{\sigma_i(\Delta \widehat{earn})},
\qquad
\Delta \widehat{earn}_{it}=\widehat{earn}_{it}-\widehat{earn}_{i,t-1}.
\]
The regression uses \(X_{i,t-1}\) (lagged one year).
\item Return controls from CRSP:
\[
ret12m_{it}=\prod_{m=1}^{12}(1+r_{i,t-m+1})-1,
\qquad
ret12m\_lag_{it}=ret12m_{i,t-1}.
\]
\end{itemize}
To avoid mechanically extreme standardized values, we require at least three non-missing first differences per firm and drop observations with absolute standardized values above 50.

\paragraph{Regression specifications.}
The six reported specifications are:
\begin{align*}
(1)\quad & Y^{pay}_{it}=a+\beta X_{i,t-1}+\varepsilon_{it},\\
(2)\quad & Y^{pay}_{it}=a+\beta X_{i,t-1}+\gamma\,ret12m\_lag_{it}+\varepsilon_{it},\\
(3)\quad & Y^{pay}_{it}=a+\beta X_{i,t-1}+\delta\,ret12m_{it}+\varepsilon_{it},\\
(4)\quad & Y^{emp}_{it}=a+\beta X_{i,t-1}+\varepsilon_{it},\\
(5)\quad & Y^{emp}_{it}=a+\beta X_{i,t-1}+\gamma\,ret12m\_lag_{it}+\varepsilon_{it},\\
(6)\quad & Y^{emp}_{it}=a+\beta X_{i,t-1}+\delta\,ret12m_{it}+\varepsilon_{it}.
\end{align*}
Standard errors are clustered at the firm (\texttt{gvkey}) level.

\paragraph{Common-sample implementation.}
The replication code optionally enforces a common sample across all six columns (same observations in payroll/workers regressions and both return controls) and optionally restricts fiscal-year ranges to align with legacy implementations.

\subsection{Turnover--Disagreement Analysis (IBES + CRSP)}\label{app:data_turnover}

\paragraph{Overview.}
The turnover analysis combines (i) a quarterly stock-market turnover measure built from CRSP, (ii) an analyst-level disagreement index built from IBES micro forecasts, and (iii) NBER business-cycle dates used to define recession indicators. The final regression sample is quarterly and spans 1982Q2--2021Q3.

\paragraph{Quarterly turnover from CRSP.}
Turnover is constructed from monthly CRSP stock data as follows:
\begin{enumerate}
\item Keep NYSE common stocks: \texttt{EXCHCD}=1 and \texttt{SHRCD} $\in \{10,11\}$.
\item For each stock-month, compute turnover as
\[
\text{turnover}_{i,t}=\frac{\text{VOL}_{i,t}}{\text{SHROUT}_{i,t}},
\]
with \texttt{SHROUT} converted to shares (\texttt{SHROUT}$\times 1000$).
\item Winsorize stock-level turnover cross-sectionally at each month at the 99th percentile.
\item Compute monthly equal-weighted and value-weighted turnover across stocks, where value weights use lagged market equity (\(|\text{PRC}_{i,t}|\times \text{SHROUT}_{i,t}\), lagged one month).
\item Aggregate monthly turnover within each quarter by summing the three monthly values and multiplying by 100:
\[
qvw_t = 100\sum_{m\in t}\text{VWTurnover}_m.
\]
\end{enumerate}
The regressions use \(qvw_t\) as the dependent variable.

\paragraph{IBES disagreement index (DI).}
Let \(e_{ijt}\) denote analyst \(j\)'s one-year-ahead earnings-growth forecast for firm \(i\) at quarter \(t\). We estimate:
\[
e_{ijt} = \alpha_{it} + \gamma_j + \beta_j g_t + \varepsilon_{ijt},
\]
where \(g_t\) is aggregate one-year-ahead realized log earnings growth (from DeLao's earnings-growth series), \(\alpha_{it}\) is a firm$\times$quarter fixed effect, and \((\gamma_j,\beta_j)\) are analyst-specific coefficients.

The index used in turnover regressions is
\[
DI_t = \text{sd}_j(\omega_{jt}), \qquad \omega_{jt}=\gamma_j+\beta_j g_t,
\]
computed across active analysts in each quarter.

\paragraph{IBES sample restrictions.}
The IBES micro sample applies:
\begin{enumerate}
\item \texttt{MEASURE}=\texttt{EPS}, \texttt{FPI}=1;
\item non-missing analyst ID, firm ID, and forecast value;
\item quarter-level winsorization of forecasts (1st/99th percentiles);
\item analyst-quarter coverage filter: at least 3 firms per analyst-quarter;
\item quarter coverage filter: at least 10 analysts in the quarter.
\end{enumerate}
These filters, together with availability of \(g_t\), imply a DI sample from 1982Q2 to 2021Q3.

\paragraph{Recession indicators and regression specifications.}
NBER peaks and troughs are used to define:
\begin{itemize}
\item \(R_t\): recession dummy;
\item \(F2_t\): dummy for the first two quarters of each recession spell.
\end{itemize}
The five turnover specifications are:
\begin{align*}
(1)\quad & qvw_t = c + \beta DI_t + u_t,\\
(2)\quad & qvw_t = c + \beta DI_t + \delta R_t + \theta (DI_t\times R_t)+u_t,\\
(3)\quad & qvw_t = c + \beta DI_t + \theta (DI_t\times R_t)+u_t,\\
(4)\quad & qvw_t = c + \beta DI_t + \delta F2_t + \theta (DI_t\times F2_t)+u_t,\\
(5)\quad & qvw_t = c + \beta DI_t + \theta (DI_t\times F2_t)+u_t.
\end{align*}
Reported standard errors are Newey--West HAC with lag 4.

\paragraph{Robustness.} As a robustness check, we control for aggregate market volatility (VIX) and market returns in the turnover regressions. 
The results are reported in Table~\ref{tab:turnover_di_robustness}.
We observe that the main results are robust to these controls: the interaction term with the recession indicator is still positive and statistically significant.

\begin{table}[!t]\centering
    \caption{Turnover Regressions with Standardized Disagreement (Robustness)}
    \label{tab:turnover_di_robustness}
    
    \resizebox{\textwidth}{!}{%
    \begin{tabular}{lccccc}
    \hline\hline
     & Baseline & Rec $\times$ DI & Rec $\times$ DI 
     & Early Rec. & Early Rec. \\
     & (1) & (2) & (3) & (4) & (5) \\
    \hline
    Disagreement (z-score) & 0.058 & 0.039 & 0.039 & 0.049 & 0.049
\\
 & (0.045) & (0.037) & (0.037) & (0.043) & (0.043)
\\
Recession &  & 0.012 &  &  & 
\\
 &  & (0.054) &  &  & 
\\
Disagreement $\times$ Recession &  & 0.174*** & 0.176*** &  & 
\\
 &  & (0.051) & (0.053) &  & 
\\
Early Rec. (first 2 qtrs) &  &  &  & -0.060 & 
\\
 &  &  &  & (0.044) & 
\\
Disagreement $\times$ Early Rec. &  &  &  & 0.127*** & 0.121**
\\
 &  &  &  & (0.046) & (0.056)
\\
VIX & 0.009*** & 0.007** & 0.007*** & 0.009*** & 0.009***
\\
 & (0.003) & (0.003) & (0.003) & (0.003) & (0.003)
\\
Market return & -0.032 & 0.008 & 0.010 & -0.042 & -0.035
\\
 & (0.141) & (0.117) & (0.114) & (0.131) & (0.131)
\\
    \hline
    Marginal effect of DI in recession 
        &  & 0.213 & 0.215 & 0.176 & 0.170 \\
        Constant & 0.180*** & 0.201*** & 0.198*** & 0.174*** & 0.179***
        \\
         & (0.057) & (0.060) & (0.053) & (0.060) & (0.060) \\
    \hline
    Observations & 127 & 127 & 127 & 127 & 127
\\
$R^2$ & 0.292 & 0.404 & 0.403 & 0.335 & 0.328
\\
Adjusted $R^2$ & 0.274 & 0.379 & 0.384 & 0.307 & 0.306 \\
    \hline\hline
    \end{tabular}%
    }
    
    \vspace{0.2em}
    {\begin{flushleft}
        \footnotesize
        Sample period: 1990Q1--2021Q3. Disagreement is standardized within sample.
        Newey--West (lag 4) standard errors in parentheses.
        * $p<0.10$, ** $p<0.05$, *** $p<0.01$.
        \end{flushleft}}
    
    
    \end{table}


\subsection{Productivity-Growth Coibion--Gorodnichenko Regressions (SPF)}\label{app:data_productivity_cg}

\paragraph{Overview.}
We construct one-quarter-ahead forecast errors and forecast revisions for two productivity measures from SPF level forecasts and realized macro data, then estimate CG regressions at horizon \(h=1\).

\paragraph{Forecast inputs from SPF.}
Using SPF mean-level forecasts for RGDP and EMP:
\[
F_t\Delta y_{t+1}=400\log\!\left(\frac{RGDP2_t}{RGDP1_t}\right),\qquad
F_t\Delta n_{t+1}=400\log\!\left(\frac{EMP2_t}{EMP1_t}\right),
\]
and similarly for two-quarters-ahead forecasts:
\[
F_t\Delta y_{t+2}=400\log\!\left(\frac{RGDP3_t}{RGDP2_t}\right),\qquad
F_t\Delta n_{t+2}=400\log\!\left(\frac{EMP3_t}{EMP2_t}\right).
\]

\paragraph{Productivity forecast measures.}
Define:
\[
F_t\Delta prod^{LP}_{t+1}=F_t\Delta y_{t+1}-F_t\Delta n_{t+1},
\]
\[
F_t\Delta prod^{TFP}_{t+1}=F_t\Delta y_{t+1}-\alpha F_t\Delta n_{t+1},\qquad \alpha=\tfrac{2}{3}.
\]
Forecast revisions are:
\[
Revision^{LP}_t = F_t\Delta prod^{LP}_{t+1} - F_{t-1}\Delta prod^{LP}_{t+1},
\]
\[
Revision^{TFP}_t = F_t\Delta prod^{TFP}_{t+1} - F_{t-1}\Delta prod^{TFP}_{t+1},
\]
with \(F_{t-1}\Delta prod_{t+1}\) obtained by lagging the \(t\)-dated two-quarters-ahead forecast.

\paragraph{Realized series.}
Realized GDP growth is constructed from real-time GDP vintages; realized employment growth is built from quarterly averages of monthly payroll employment (PAYEMS):
\[
\Delta y_{t+1}=400\log\!\left(\frac{RGDP_{t+1}}{RGDP_t}\right),\qquad
\Delta n_{t+1}=400\log\!\left(\frac{PAYEMS_{t+1}}{PAYEMS_t}\right).
\]
Hence:
\[
\Delta prod^{LP}_{t+1}=\Delta y_{t+1}-\Delta n_{t+1},\qquad
\Delta prod^{TFP}_{t+1}=\Delta y_{t+1}-\alpha\Delta n_{t+1}.
\]

\paragraph{CG regressions.}
For each measure \(m\in\{LP,TFP\}\), estimate:
\[
FE^m_t = c + \beta_{CG}^m\,Revision^m_t + u_t,
\]
where \(FE^m_t=\Delta prod^m_{t+1}-F_t\Delta prod^m_{t+1}\).
Reported standard errors are Newey--West HAC with lag 4.

\paragraph{Sample.}
The productivity CG table uses 2004Q1--2025Q2 (86 quarterly observations in each column).

\subsection{Construction of Table 5 Data Moments and Bootstrap Intervals}\label{app:data_table5_moments}

\paragraph{Overview.}
The data targets in Table 5 are constructed by ingesting the raw series described below and computing the moments and bootstrap uncertainty intervals reported in the table.
All moments are reported in annualized percentage points.

\paragraph{Data sources (key series).}
For the moments requested in Table 5, we use:
\begin{itemize}
\item \textit{Excess returns and dividend growth}: Shiller historical market dataset (CSV mirror), using real equity returns and nominal long rates.
\item \textit{Consumption growth}: BEA/FRED annual real nondurables and services series, \texttt{DNDGRG3A086NBEA} and \texttt{DSERRG3A086NBEA}, divided by annual population \texttt{B230RC0A052NBEA}.
\end{itemize}

\paragraph{Series and definitions.}
We construct five data objects:
\begin{enumerate}
\item \textit{Real interest rate}: annual ex-post real long rate from Shiller,
\[
r^{real}_t=r^{nominal}_t-\pi_t,
\]
where \(r^{nominal}_t\) is the Shiller long rate and \(\pi_t\) is realized annual inflation computed as a log change in the Shiller CPI (December-to-December, in our baseline construction).
\item \textit{Excess equity return}: annual real equity return minus the real interest rate.
\item \textit{Consumption growth}: annual log growth of real per-capita nondurables plus services consumption,
\[
\Delta c_t = 100\Delta \log\!\left(\frac{DNDGRG3A086NBEA_t+DSERRG3A086NBEA_t}{B230RC0A052NBEA_t}\right).
\]
\item \textit{Dividend growth}: annual log growth of annual-average real dividends from Shiller.
\item \textit{Log-hours volatility}: standard deviation of HP-filtered (\(\lambda=100\)) annual log total hours, where annual hours are built from CPS/BLS series as
\[
hours_t=\frac{LNU02005053_t \times LNU02005054_t}{LNU00000000_t-LNU00000097_t}.
\]
\noindent We access these monthly CPS series via the BLS API, construct \(hours_t\) at the monthly frequency, and then take annual averages before applying the HP filter to the annual log series. This construction mirrors \citet{rogerson2011search}.
\end{enumerate}

\paragraph{Baseline sample windows used for Table 5.}
The replication uses:
\begin{itemize}
\item financial and dividend moments: 1890--2025;
\item consumption growth: 1945--2025 (paper-table definition above);
\item hours volatility: 1976--2025.
\end{itemize}

\paragraph{Annualization and reported moments.}
For each variable \(x_t\), we report
\[
\widehat{\mu}_x = \text{mean}(x_t), \qquad
\widehat{\sigma}_x = \text{sd}(x_t),
\]
already scaled in annual percentage points by construction.

\paragraph{Bootstrap standard errors and intervals.}
To report uncertainty in the data column, we compute moving-block bootstrap standard errors for both \(\widehat{\mu}_x\) and \(\widehat{\sigma}_x\):
\begin{itemize}
\item number of bootstrap replications: \(B=1000\);
\item seed: 12345;
\item block length: \(l=\lfloor n^{1/3}\rfloor\) (unless manually overridden), where \(n\) is the sample size of the variable.
\end{itemize}
For each replication \(b\), we resample contiguous blocks of annual observations (years) with replacement, form a pseudo-sample \(x_t^{(b)}\), and compute \(\widehat{\mu}_x^{(b)}\) and \(\widehat{\sigma}_x^{(b)}\). Standard errors are:
\[
\widehat{se}(\widehat{\mu}_x)=sd_b\!\left(\widehat{\mu}_x^{(b)}\right), \qquad
\widehat{se}(\widehat{\sigma}_x)=sd_b\!\left(\widehat{\sigma}_x^{(b)}\right).
\]
Table 5 reports \(\pm 2\) s.e. intervals:
\[
\left[\widehat{\mu}_x-2\widehat{se}(\widehat{\mu}_x),\,\widehat{\mu}_x+2\widehat{se}(\widehat{\mu}_x)\right], \qquad
\left[\widehat{\sigma}_x-2\widehat{se}(\widehat{\sigma}_x),\,\widehat{\sigma}_x+2\widehat{se}(\widehat{\sigma}_x)\right].
\]
